% C.43 candidate for canonical integration after final root read.
\begin{statement}{C.43}{Proved}{Exactly two binary classes remain above the smaller finite quotient}
Use the actual finite source and normal quotients of C.42.  Put
\[
 G=G_{176},\qquad Q=G_{104},\qquad
 A=\ker(G\to Q)\simeq(C_2)^{72}.
\]
The actual finite extension $1\to A\to G\to Q\to1$ splits.
The natural quotient map to $G_{86}$ induces a split exact sequence
\[
 0\longrightarrow(C_2)^2\longrightarrow H_2(G,\mathbb Z_2)
   \longrightarrow H_2(G_{86},\mathbb Z_2)\longrightarrow0.
 \tag{C.43.1}
Its kernel is the actual image of
$H_0(Q,\Lambda^2A)$ and has exactly the displayed order.
The splitting depends on a choice of group section; the natural map
and its kernel do not.

Consequently a verified two-primary lifting calculation in $G_{86}$
would leave two binary values to determine in the original source.
Neither value is computed by this statement.  In particular C.42
and (C.43.1) do not yet produce a full iterated-descent lift.
\end{statement}
\begin{proof}
The original pc generators $g_{105},\ldots,g_{176}$ give the stated
basis of $A$.  The actual normal subgroup
\[
       N=\ker(Q\to G_{86})=\langle g_{87},\ldots,g_{104}\rangle
                            \simeq(C_3)^{18}
 \tag{C.43.2}
has odd order.  These normality and elementary abelian assertions
are checked from every original relation, as in C.42.

\paragraph{An exact rank certificate on the original module.}
For every one of the $104$ quotient generators, read its action on
$A$ directly from the original conjugation words.  The matrices
are invertible over $\mathbb F_2$.  The stacked rows $M_g-I$ for
the $18$ generators in (C.43.2) have rank $72$, proving
\[
                              A_N=0.
 \tag{C.43.3}
The exterior square has dimension $\binom{72}{2}=2556$.
Stacking the exterior action minus the identity for all $104$
quotient generators gives $265\,824$ actual rows and rank $2554$.
Hence
\[
                       (\Lambda^2A)_Q\simeq(C_2)^2.
 \tag{C.43.4}

Both rank inequalities in (C.43.4) have small exact certificates.
The full program retains $2554$ independent original action rows.
A separate checker reconstructs them from the original presentation
and eliminates them with lowest-coordinate pivots, whereas the first
calculation used highest-coordinate pivots.  This independently
proves rank at least $2554$.  The checker also reconstructs two
explicit independent alternating forms on $A$ and verifies their
invariance under every one of the $104$ original matrices.  They
give rank at most $2554$.  Thus no unprocessed part of the full
row system is presumed to vanish.

The full data are in
\texttt{pbc\_two\_top\_module.json}; the fast independent
certificate is
\texttt{pbc\_two\_top\_certificate.py} and its retained output.
The second checker rederives all action matrices from the source
hash of C.42, rather than trusting only the first program's rank.
The independent audit also checks the literal normal suffixes and
their elementary kernels in the original group.

\paragraph{The actual extension splits.}
Because $|N|$ is odd, averaging is defined on the exponent-two
module $A$.  It identifies invariants with coinvariants and makes
both functors exact.  Equation (C.43.3) therefore gives
\[
 A^N=0,\qquad H_i(N,A)=H^i(N,A)=0\quad(i>0).
\]
Apply both homology and cohomology Hochschild--Serre to the actual
normal subgroup $N$ in $Q$.  Every second-page coefficient vanishes,
including degree zero, so
\[
                    H_i(Q,A)=H^i(Q,A)=0\qquad(i\ge0).
 \tag{C.43.5}
In particular $H^2(Q,A)=0$.  Classification of extensions by this
cohomology group gives an actual homomorphism section $Q\to G$
\cite[II, Example~1.18(b), Remark~1.35(b)]{PBC_MilneCFT}.
This proves the group-level splitting; no integral Lie-algebra or
associated-graded splitting is substituted for it.

\paragraph{All possible differentials in degree two.}
For $A$ normal in $G$, use the homology spectral sequence
\[
 E^2_{a,b}=H_a(Q,H_b(A,\mathbb Z_2))
                              \ \Longrightarrow H_{a+b}(G,\mathbb Z_2).
 \tag{C.43.6}
As the groups are finite, these coefficients retain exactly the
two-primary part of their integral homology.  The integral identity
$H_2(A,\mathbb Z)=\Lambda^2A$ identifies the left edge in total
degree two.  By (C.43.4)--(C.43.5),
\[
 E^2_{0,2}=(C_2)^2,\qquad E^2_{1,1}=0,\qquad
                         E^2_{2,0}=H_2(Q,\mathbb Z_2).
 \tag{C.43.7}
There are only two possible incoming differentials to $E_{0,2}$.
The $d_2$ source is $H_2(Q,A)=0$.  The $d_3$ source is in the
bottom row $E_{3,0}$, and this differential is zero because the
actual section makes every bottom-row class permanent: projection
$H_i(G,\mathbb Z_2)\to H_i(Q,\mathbb Z_2)$ is split surjective.
There are no later incoming differentials and no outgoing
differentials from the left edge.

The same section protects $E_{2,0}$.  Thus the filtration gives
\[
 0\longrightarrow(C_2)^2\longrightarrow H_2(G,\mathbb Z_2)
        \longrightarrow H_2(Q,\mathbb Z_2)\longrightarrow0,
\]
and the section splits this exact sequence as well.  This proves
that the kernel is exactly $(C_2)^2$, rather than merely a possible
quotient of that group.  Finally $Q\to G_{86}$ has the odd-order
kernel $N$, so its homology spectral sequence gives a natural
isomorphism on $H_2(-,\mathbb Z_2)$.  Composing proves (C.43.1).

\paragraph{What the finite calculation does not decide.}
The two explicit invariant alternating forms both vanish on
$a\wedge s(a)$ for every basis vector $a$ of $A$, where
$s=g_2$ is the actual involution.  The map
$a\mapsto B(a,s(a))$ is additive: its cross terms cancel because
$s$ is an involution preserving $B$.  Thus these values vanish for
every $a\in A$.  All norms $a+s(a)$ lie in the span of the $36$
retained top-kernel centralizer generators.  This particular norm
calculation therefore does not provide the two missing values.

A separate modular search in $G_{86}$ has been retained, including
its complete proof records.  Its first independent original-source
replay rejects the raw target normal form.  It supplies no verified
quotient identity and is not used anywhere in this proof.  Even a
successful quotient replay would still require the two values in
(C.43.1) before a conclusion about the original finite obstruction.
\end{proof}
